Paragraph 1 theme: Blinking as a behavioural indicator of fatigue driving.
Blinking is a physiological process that involves the closure and reopening of the eyelids. Its primary functions are to protect the ocular surface, maintain corneal moisture, and support visual comfort. Beyond these physiological functions, blink behaviour has been increasingly recognised as an important indicator in fatigue-driving research. Variations in blink-related features, such as blink frequency, blink duration, eye-closure duration, and blink patterns, can reflect changes in alertness and cognitive state. As fatigue increases, drivers may exhibit slower or prolonged blinks, which can indicate reduced vigilance and a higher risk of impaired driving performance. Therefore, blink behaviour provides a useful and non-invasive marker for monitoring driver fatigue.